\subsection{Auch die Skalen besitzen einen gemeinsamen Variationsansatz}
Im Metrologieoriginal ist die Planckfrage stärker formuliert als die bloße Wahl einer Einheit. Mit der ersten positiven Nahtfrequenz $\lambda_\Sigma$ und $\rho=\chi^2/\lambda_\Sigma^2$ lautet das Potential in Rand-Einheiten
\begin{equation}
\widehat V(\rho)=-6\widehat{\mathcal E}_D\rho
+\widehat\Gamma_D^{(4)}(\rho)
+\rho^2[-\log\det_{\rm adm}(I-\widehat U_\Sigma(\rho))].
\end{equation}
Unter $\widehat V'(0^+)<0$, $\widehat V''(\rho)>0$ für $\rho>0$ und $\widehat V'(\rho)\to+\infty$ gibt es genau ein positives stationäres $\rho_\star$. Dann gilt
\begin{equation}
\frac{\bar M_{\rm Pl}^2}{\lambda_\Sigma^2}=\frac{\rho_\star}{2\pi^2},\qquad
G_N\lambda_\Sigma^2=\frac{\pi}{4\rho_\star}.
\end{equation}
Das ist ein wirklicher bedingter Eindeutigkeitssatz. Sein Zahlenwert verlangt den konkreten Operator und die Prüfung der genannten Voraussetzungen. Eine Konvexitätsannahme allein berechnet $\rho_\star$ noch nicht.

Der ursprüngliche EW-/Vakuumenergieanschluss liefert, mit $A=\alpha^{-1}$,
\begin{equation}
\frac{v_{\rm EW}}{\bar M_{\rm Pl}}=5\beta_{\rm rad}^2e^{-(A+\delta_{\rm ph})/5},
\qquad \frac{\rho_\Lambda}{\bar M_{\rm Pl}^4}=\frac3{4\pi^2}e^{-2A}.
\end{equation}
Elimination derselben Größe $A$ ergibt exakt
\begin{equation}\label{eq:physicalscalebridge}
\boxed{\frac{\rho_\Lambda}{\bar M_{\rm Pl}^4}
=\frac3{4\pi^2}e^{2\delta_{\rm ph}}
\left(\frac{v_{\rm EW}}{5\beta_{\rm rad}^2\bar M_{\rm Pl}}\right)^{10}.}
\end{equation}
Die Branchwahl, Higgsnormalisierung und der Transportpol bleiben die ausdrücklich zu identifizierenden Eingänge. Innerhalb dieses Vertrags sind die beiden Skalen nicht unabhängig einstellbar.

Die Scalaronroute benutzt zusätzlich $M_{\rm scal}/\bar M_{\rm Pl}=c_3^{7/2}$ und liefert bedingt $A_s=N_*^2c_3^7/(24\pi^2)$, $n_s=1-2/N_*$, $r=12/N_*^2$. $N_*$ bleibt Szenarioeingang. Bei ihrer stärksten Normierungsbegründung ist eine konkrete Prüfung offen: Der in v36 isolierte $R^2$-Summand ist nicht der vollständige $a_4$ des dort verwendeten gewöhnlichen Spin-Diracoperators. Dessen vollständige gravitative Kombination ist proportional zu
\begin{equation}
5R^2-8R_{\mu\nu}^2-7R_{\mu\nu\rho\sigma}^2=11E_4-18C^2,
\end{equation}
modulo Divergenzen. Der relative, projizierte oder torsionshaltige TFPT-Masteroperator muss daher seinen vollständigen Koeffizienten liefern; die Standardrechnung allein liefert die behauptete Starobinsky-Normierung nicht \cite{chamseddineconnes}. Das grenzt diesen konkreten Begründungsweg ein und verwirft nicht die anderen Skalenrelationen.

\subsection{Der gemeinsame Operator ist im Original bereits präzise vorgegeben}
Der Flavoransatz schreibt tatsächlich
\begin{equation}
P_{\rm str}=P_{\Sigma,+}P_{\rm sing}P_\Theta,\quad
K_f^{\rm str}=(P_{\rm str}D_f^\dagger D_fP_{\rm str})^{-1/2},\quad
Y_f=\operatorname{Hol}(\nabla_F^\star)K_f^{\rm str},
\end{equation}
auf dem zulässigen, invertierbaren Projektionsraum. Der Transportpol wird über das ursprüngliche Determinantenfunktional von $D_{y,\delta}=yI-\delta U_6$ bestimmt. Für die neue Quelle ist deshalb die konkrete Gleichheit $Y_f^{\rm source}=Y_f$ mit denselben Familienlegs, $K/L$-Gradierungen und derselben Norm zu zeigen. Die neue E8-Quelle enthält nicht automatisch die Flavorhierarchie, nur weil sie die richtige Stromalgebra besitzt.

Zugleich muss ihre primitive Nahtauslesung das oben angegebene $U_\Sigma^{\min}$ und die $\alpha$-Variation ergeben; die Skalenvariation muss das genannte Planckpotential erzeugen. Die ursprüngliche Forderung nach einer gemeinsamen relativen Determinantenfamilie verbindet damit bereits vorhandene konkrete Gesetze. Sie ist wesentlich enger als ``finde irgendein Funktional''.

\key{\textbf{Veränderte Arbeitspriorität:} Die fundamentale Quelle muss am ursprünglichen markierten Flavor- und Nahtoperator identifiziert werden. $M_F$, $Q,\Sigma,K,L$, die $48/41$-Gewichte und die Skalenvariationen sind dabei gleichzeitig zu erhalten. Weitere Rechnungen nur an der neuen symmetrischen E8-Quelle ersetzen diesen Anschluss nicht.}

\subsection{Neutrinos: Die vorhandenen Herleitungen gemeinsam lösen}\label{sec:neutrinojoint}

Der Originalbestand liefert einen endlichen Diracoperator mit Dirac- und Majoranablock, den tatsächlichen Seesaw-Ausdruck und konkrete PMNS-Kanäle. Auf dem ursprünglichen Zweig stehen unter anderem $\sin^2\theta_{12}=1/3-\phi_0/2$, $\sin^2\theta_{13}=\phi_0e^{-5/6}$ und die hexagonale Phasenlesart $\delta_{\rm PMNS}=240^\circ$. Der spätere Pentagonzweig aus \file{FLAV.NUSCALE.06} benutzt dagegen eine andere Mischung und Phase; er ist ein ausdrücklich bedingter Alternativkandidat. Beide Zweige dürfen nicht zu einer einzigen voraussetzungslosen Herleitung zusammengezogen werden.

Die Scalaronverbindung ist eine echte sektorübergreifende Bedingung:
\begin{equation}
M_3=3M_{\rm scal},\qquad
M_{\rm scal}=c_3^{7/2}\bar M_{\rm Pl},\qquad y_{\nu,3}=y_t
\quad\text{im erklärten minimalen Matching.}
\end{equation}
Der Original-Ein-Schleifenlauf ergibt damit $m_{3,0}=0.05124031018\,\mathrm{eV}$. Die gemeinsame Einbettung dieses Matchings, der vollständigen Mischung und der Materieasymmetrie muss jedoch denselben Yukawaoperator verwenden.

Das wird am eingefrorenen späteren Kandidaten besonders konkret. Er setzt
\begin{equation}
M_R=M_{\rm scal}\diag(\epsilon,2\epsilon,3),\quad
\epsilon=\frac{\phi_0^2}{10},\qquad
Y_\nu=\diag(0,y_2,y_3)U^\dagger.
\end{equation}
In dieser Konvention tragen die Zeilen die schweren Neutrinos. Für jedes unitäre $U$ ist dann
\begin{equation}
Y_\nu Y_\nu^\dagger=\diag(0,|y_2|^2,|y_3|^2).
\end{equation}
Das leichteste schwere Neutrino ist entkoppelt. Sein Washout und seine Zerfallsquelle verschwinden. Der separat erfolgreiche Boltzmannlauf verwendet dagegen $\widetilde m_1=m_3/10$ sowie einen CP-Quellterm aus einer Davidson--Ibarra-Schranke mal $|\sin240^\circ|$. Diese beiden Teilrechnungen besitzen daher noch keinen gemeinsamen Yukawaoperator. Die Schranke ist nicht automatisch eine Gleichheit und die PMNS-Phase nicht automatisch die schwere Zerfallsphase \cite{davidsonibarra}.

\textbf{Die Fortsetzung benutzt einen bereits vorhandenen Originalansatz.} Die Verträge \file{NU.CI.01/02} beschränken die allgemeine Casas--Ibarra-Zuordnung \cite{casasibarra} durch Familienzyklus, Nahtreflexion, komplexe Orthogonalität und Determinante eins. Zur Unterscheidung vom Flavorresiduum $R$ heißen ihre Matrix und ihr Fouriertransformator hier $R_{\rm CI}$ und $F_3$:
\begin{align}
R_{\rm CI}(z)&=F_3\diag(1,e^z,e^{-z})F_3^\dagger,
\quad (F_3)_{jk}=\frac{\omega^{jk}}{\sqrt3},\quad \omega=e^{2\pi i/3},\\
R_{\rm CI}^{\mathsf T}R_{\rm CI}&=I,\qquad z\in\R.
\end{align}
Dies ist der identitätsverbundene reelle Zweig. Die zweite reelle Komponente mit beiden nichttrivialen Fourier-Eigenwerten negativ wird unten durch die Washoutbedingung ausgeschlossen. Die Symmetriebedingungen allein wählen noch keinen Wert von $z$.

Mit $v_D=v_{\rm EW}/\sqrt2$ lautet die gemeinsame Matrix am selben Matchingmaßstab
\begin{equation}
Y_\nu=\frac{i}{v_D}\sqrt{M_R}\,R_{\rm CI}(z)\sqrt{m}\,U^\dagger,
\qquad
-v_D^2Y_\nu^{\mathsf T}M_R^{-1}Y_\nu=U^*mU^\dagger.
\end{equation}
Das setzt ein gegebenes Niederenergiespektrum korrekt fort; es behauptet noch keine unabhängige Auswahl dieses Spektrums aus dem Ursprung. Im summierten schweren Kanal fällt $U$ aus $Y_\nu Y_\nu^\dagger$ heraus. Bei einzelnen Leptonflavors gilt diese Vereinfachung nicht.

Schreibe $c=\cosh z$, $h=\sinh z$. Die erste Zeile ist $(d,a-ib,a+ib)$ mit
\begin{equation}
d=\frac{1+2c}{3},\qquad a=\frac{1-c}{3},\qquad b=\frac{h}{\sqrt3}.
\end{equation}
Für $m_1=0$ erhält man exakt
\begin{equation}
\widetilde m_1=(m_2+m_3)f(z),\qquad
f(z)=\frac{(c-1)^2+3h^2}{9}.
\end{equation}
Der bestehende bedingte Anker $\widetilde m_1=m_3/10$ legt daher den Betrag fest:
\begin{equation}\label{eq:nuactualselection}
\boxed{\cosh z_\star=
\frac{1+3\sqrt{1+4m_3/[10(m_2+m_3)]}}4.}
\end{equation}
Es bleiben genau die beiden Orientierungen $z=\pm z_\star$. Die zweite reelle Komponente hat mindestens $\widetilde m_1=4(m_2+m_3)/9$ und kann den Anker nicht erfüllen. Der zuvor hypothetisch eingesetzte Selektor $z=6\log(3/2)$ ist damit tatsächlich prüfbar: Für den eingefrorenen Massensatz liefert er $0.7871\,\mathrm{eV}$ Washout statt $0.005124\,\mathrm{eV}$. Diese zwei bestehenden Bedingungen passen in dieser Klasse nicht zusammen.

\textbf{Washout und CP sind jetzt gemeinsam gekoppelt.} Der tatsächliche CP-Zähler lautet
\begin{equation}
\operatorname{Im}\sum_jm_j^2(R_{\rm CI})_{1j}^2
=2ab(m_3^2-m_2^2).
\end{equation}
Daraus folgt in der hierarchischen ungeflavorten SM-Näherung
\begin{equation}
\varepsilon_1^{\rm hier}
=\frac{3M_1(m_3-m_2)}{16\pi v_D^2}\,
\frac{\sqrt3\sinh z}{2\cosh z+1},
\end{equation}
mit einheitlichen Masseneinheiten und der angegebenen CP-Vorzeichenkonvention. Die gleiche Norm, die den Washout festlegt, bestimmt damit den CP-Betrag. Besonders kompakt lässt sich $z$ vollständig eliminieren. Setze
\begin{equation}
w=\frac{\widetilde m_1}{m_2+m_3},\qquad s=\sqrt{1+4w},\qquad
\mathcal C=\frac{16\pi v_D^2\varepsilon_1^{\rm hier}}{3M_1(m_3-m_2)}.
\end{equation}
Dann gilt innerhalb dieser ursprünglichen Klasse exakt
\begin{equation}\label{eq:nuwashoutcp}
\boxed{\mathcal C^2=\frac{w(3s+5)}{(s+1)^3}.}
\end{equation}
Für endliches $w$ bleibt $|\mathcal C|<\sqrt3/2$. Der früher eingesetzte Betrag $|\sin240^\circ|=\sqrt3/2$ wird in dieser Familie erst bei $|z|\to\infty$ erreicht, während zugleich der Washout divergiert. Er ist hier eine Grenzamplitude, keine unabhängig einsetzbare endliche Phase. Diese Schlussfolgerung benutzt die vorhandene Familie und den vorhandenen Anker; sie führt keine neue angepasste Textur ein.

\textbf{Die Rückwirkung auf die Masse wird mitgelöst.} Würde man den alten Wert $m_{3,0}$ unverändert behalten, ergäbe die Washoutgleichung $|z|=0.4829422092$. Dann stiege aber die dritte schwere Yukawa-Zeilennorm um den Faktor $1.08594$; insbesondere wäre auch der größte Singularwert größer als der alte Top-Matchingwert. Deshalb darf der alte Massennahetreffer nicht zusammen mit dieser neuen Matrix als unverändert bestätigt ausgegeben werden.

Die ursprüngliche Einfamilienrechnung definiert keine vollständige $3\times3$-Basiszuordnung. Als notwendige basisinvariante Fortsetzung ihres dominanten Yukawa-Matchings verlangen wir
\begin{equation}
\sigma_{\max}(Y_\nu)=y_t(M_3).
\end{equation}
Dies ist ausdrücklich schwächer als eine vollständig begründete Matrixgleichheit $Y_\nu=Y_u$ im minimalen $10_H$-Modell. Mit der im Original verwendeten gemeinsamen Rückskalierung $R_\downarrow=0.7944014077$ kürzt sich deren gemeinsamer Faktor aus der relativen Normgleichung heraus. Setze
\begin{equation}
S=\diag\bigl(\sqrt{M_1/M_3},\sqrt{M_2/M_3},1\bigr),
\qquad z=z(m)\text{ aus Gleichung~\eqref{eq:nuactualselection}}.
\end{equation}
Dann lautet das gleichzeitig zu lösende System kompakt
\begin{equation}\label{eq:nujointmassnorm}
\boxed{\lambda_{\max}\!\left[
S R_{\rm CI}(z(m))\diag(0,m_2,m)R_{\rm CI}(z(m))^\dagger S
\right]=m_{3,0}.}
\end{equation}
Die nichtverschwindenden Eigenwerte stammen aus einer $2\times2$-Gramform. Auf dem positiven Zweig wachsen deren Diagonaleinträge und der Betrag ihres Nebendiagonaleintrags streng mit $m$. Ihr größter Eigenwert wächst daher streng, beginnt bei $(M_2/M_3)m_2<m_{3,0}$ und geht gegen unendlich. Es gibt genau eine positive Massenlösung; die beiden Orientierungen bleiben konjugiert.

Mit den unveränderten schweren Massen aus der eingefrorenen v3-Datei und deren $m_2=0.008654478609\,\mathrm{eV}$ folgt
\begin{equation}
\boxed{m_3=0.0435180861582\,\mathrm{eV},\quad
z=\pm0.477382171162,\quad
\widetilde m_1=0.00435180861582\,\mathrm{eV}.}
\end{equation}
Dabei ist $m_2$ im Original ausdrücklich daten-kalibriert. Washout-Anker, CI-Symmetrieklasse und Scalaron-Matching bleiben bedingt. Dies ist eine eindeutige gemeinsame Konsistenzlösung der angegebenen notwendigen Bedingungen, keine neue unabhängige Neutrinovorhersage und kein voller $10_H$-Abschluss. Die vollständige $3\times3$-Rechnung und die reduzierte Gramrechnung wurden unabhängig gegengeprüft.

Im vorhandenen Boltzmannpfad wurde danach die tatsächliche CP-Quelle dieser Matrix verwendet. Unter der bisherigen ungeflavorten $N_1$-Näherung ergibt sich:
\begin{center}
\begin{tabular}{lrr}
\toprule
 & Hierarchische Schleife & Endliche Massenschleifen\\\midrule
$|\varepsilon_1|$ & $1.57943\cdot10^{-7}$ & $1.36236\cdot10^{-7}$\\
$\kappa_f$ & $0.1129051$ & $0.1129051$\\
$|\eta_B|$ & $1.71193\cdot10^{-10}$ & $1.47665\cdot10^{-10}$\\\bottomrule
\end{tabular}
\end{center}
Die endlichen, nichtresonanten Ein-Schleifenfunktionen \cite{covirouletvissani} sind hier wichtig, weil $M_2=2M_1$ keine starke Hierarchie ist. Der in dem alten Programm gespeicherte Vergleichswert beträgt $6.1\cdot10^{-10}$; der gemeinsame Anschluss erreicht in dessen Näherung rund $24\,\%$ davon. Dieser Vergleich benutzt den eingefrorenen Datenstand, keine heute neu recherchierte Messung.

Die Rechnung erfolgte direkt mit \file{kappa_f_ode} im bestehenden \file{experiments/ftransfer/leptogenesis_boltzmann/fboltzmann_solve.py}. Die ursprünglichen schweren Massen wurden aus \file{experiments/nu-scalaron-falsification/hypotheses/nu_scalaron_v3.yaml} übernommen. Die Algebra benutzt \file{nu01_casas_ibarra_seam.py} und prüft die Hypothese aus \file{nu02_transport_selector.py} im Verzeichnis \file{experiments/theory-contracts/}. Eine vollständige Flavor-Dichtematrixrechnung mit mehreren schweren Spezies und Schwellen-RG wurde nicht ausgeführt. Die ODE übernimmt bewusst die bisherige Niedrigenergie-/Matching-Näherung des Originals. Insbesondere beweist das verlorene Nahetreffen bei dessen großer RG-Schemenunsicherheit noch keinen Ausschluss der gesamten Leptogeneseroute. Auch die physische Zuordnung des P1-Orientierungszeichens zu positivem Baryonenüberschuss ist dadurch nicht hergeleitet.


\picturetext{Für dasselbe schwere Neutrino sind Kopplungsstärke, erzeugte Asymmetrie und ihr späteres Auswaschen miteinander verbunden. Man kann diese drei Regler nicht einzeln aus drei erfolgreichen Rechnungen übernehmen. Die gemeinsame Gleichung zeigt nun genau, welche Einstellungen zugleich möglich sind.}

\textbf{Bedeutung für die fundamentale Frage:} Die vorhandene Neutrinogeometrie liefert bereits den nötigen nichttrivialen Kanal. Ihre gemeinsamen Bedingungen lösen einen Teil der bisher freien Zuordnung und schließen eine konkrete ältere Gap-Identifikation aus. Die ursprünglichen Massen-, Mischungs- und Kosmologieansätze sind dadurch engere Anforderungen an dieselbe Quelle geworden. Der Befund betrifft diesen gekoppelten Ansatz und diese Transfernäherung; er widerlegt weder den gesamten Neutrinozweig noch den unabhängigen kosmologischen $\Omega_b$-Readout.



\subsection{Nachweise der erneuten Prüfung}
Die Originalkennungen sind insbesondere \file{FLAV.LEPTONC.01}, \file{FLAV.RIGID.01/02}, \file{FLAV.CKM.01}, \file{FLAV.PMNS.02}, \file{FLAV.NUSCALE.05/06}, \file{EM.WARD.01} und \file{ALPHA.QUILLEN.EXACT.01}. Neu ausgeführt wurden v20 (12/12), v71 (7/7), v118 (6/6), v88 (11/11), v3 (14/14) und v485 (5/5); diese Läufe prüfen ihre angegebenen Formeln und Operatorrechnungen. Eine neue Ausführung der gesamten Lean- oder Repository-Suite wird damit nicht behauptet.

Maßgebliche Vertiefungen: \file{tfpt_2_standard_model.tex}, Zeilen 131, 310, 450, 573, 1751 und 1862; \file{tfpt_1_architecture_e8.tex}, Zeilen 3968--4050; \file{_archive/tfpt-45/source_extracts/03_em_flavor_source.tex}, Zeilen 94--230, 811--846 und 1257--1273; \file{05_metrology_source.tex} im selben Ordner, Zeilen 622 und 1370--1525. Das spätere \file{source-flavor-joint-dictionary-20260921/PROOF.md} erhält die vollständigen markierten Operatorprodukte. Sein endliches Wörterbuch ersetzt noch nicht die Auswahl ihrer physischen Quellenmetrik. Der ältere Archiv-Monodromiesatz wird hier nicht für die vollständige Auswahl verwendet: Seine drei angegebenen konjugiert-diagonalen Monodromien multiplizieren sich zu $I$, wodurch die vierte die gleichzeitig verlangte nichttriviale Punkturspektralbedingung verletzt. Die ausführbare v117/v118-Familienmonodromie bleibt die geprüfte Grundlage der hier erhaltenen Leptonformeln.

\section{Fortsetzung: Der gemeinsame Flavor- und Nahtanschluss}\label{sec:operatorbridge}

\key{\textbf{Die gemeinsame Verbindung ist konkreter geworden.}
Das ursprüngliche Fuchssystem liefert jetzt einen wirklichen Wegoperator, der exakt zum bereits verwendeten Sechsertransport äquivalent ist. Auf der anderen Seite liefert die ursprüngliche E8-Quelle eine konkrete neutrale Anregung mit Gewicht zwei und Normquadrat 48. Die gemeinsame Aufgabe ist dadurch keine Suche nach irgendeiner passenden Matrix mehr: Es ist die Herleitung der angegebenen Nahtkopplung und ihrer gemeinsamen Antwort mit den ursprünglichen Flavor-Defekten.}

\subsection{Der ganze endliche Operator besitzt bereits ein gemeinsames Wörterbuch}
Der Originalvertrag \file{source-native-clock-flavor-bridge-20260922/PROOF.md} liefert eine feste invertierbare Matrix $S$ mit
\begin{equation}
S^\dagger S=N,\qquad SPS^{-1}=M,\qquad
S(-G)S^{-1}=U,\qquad Q_c=SQS^{-1}.
\end{equation}
$N$ ist die dort ausgewählte positive Familienform, $P$ der ganzzahlige Dreierzyklus, $M$ die Familienmonodromie und $U$ die Vierteldrehung. Die Orientierungsgerade ist in $-G$ ausdrücklich mitgeführt. Für den verdoppelten Raum folgt deshalb
\begin{equation}\label{eq:fullflavorintertwiner}
(S\oplus S)\bigl[yI-\delta(P\oplus(-P))\bigr](S\oplus S)^{-1}
=yI-\delta(M\oplus(-M)).
\end{equation}
Mit dem zu $N$ gehörenden Adjungieren werden gleichzeitig $D^\dagger D$, Resolventen, inverse Quadratwurzeln, Polarfaktoren und Determinanten transportiert. Das ist stärker als die Gleichheit von Eigenwertlisten. Es ist keine neue freie Wahl eines Flavoroperators.

Die ursprüngliche positive Transportantwort und der inverse Operator bleiben verschieden:
\begin{equation}
D=V|D|,\qquad K=|D|^{-1},\qquad
(D^\dagger D)^{-1}=K^2,\qquad D^{-1}=KV^\dagger .
\end{equation}
Die zusätzlich im Yukawagesetz benutzte Familienholonomie ist dadurch noch nicht als $V^\dagger$ identifiziert. Ebenso ist die ursprüngliche Wortgewichtung $\lambda_Y^L$ kein automatischer Ersatz der geometrischen Reihe in $\delta/y$. Beide tragen eigene, bereits markierte Daten.

\subsection{Die Vierer- und Sechsergeometrie gehören zu demselben Überlagerungsproblem}\label{sec:flavorcover}
\textbf{Zunächst lässt sich die bisher numerische Fuchs-Monodromie analytisch schließen.}
Das ursprüngliche System aus v117 ist $\Psi'=A(z)\Psi$, mit
\begin{equation}
A(z)=\sum_{k=0}^3\frac{U^kA_0U^{-k}}{z-i^k},\qquad
A_0=\begin{pmatrix}
1/2&\sqrt2/6&0\\
\sqrt2/6&1/4&\sqrt5/12\\
0&\sqrt5/12&1/4
\end{pmatrix},\qquad U=\diag(1,i,-i).
\end{equation}
Die Substitution $x=z^4$, $\Psi(z)=\diag(1,z,z^3)v(x)$ liefert exakt
\begin{equation}
\frac{dv}{dx}=\left(\frac{B_0}{x}+\frac{B_1}{x-1}\right)v,\qquad
B_0=\begin{pmatrix}
0&0&0\\-\sqrt2/6&-1/4&0\\0&-\sqrt5/12&-3/4
\end{pmatrix},\quad B_1=A_0,
\end{equation}
und
\begin{equation}
B_\infty=-B_0-B_1
=\begin{pmatrix}
-1/2&-\sqrt2/6&0\\0&0&-\sqrt5/12\\0&0&1/2
\end{pmatrix}.
\end{equation}
Die lokalen Exponenten bei $0,1,\infty$ sind somit
$(0,-1/4,-3/4)$, $(0,1/3,2/3)$ und $(-1/2,0,1/2)$.

Bei $\infty$ muss die ganzzahlige Exponentendifferenz geprüft werden; die Eigenwerte allein genügen dort nicht. Mit
\begin{equation}
S_\infty=\begin{pmatrix}
1&-\sqrt2/3&\sqrt{10}/36\\
0&1&-\sqrt5/6\\0&0&1
\end{pmatrix},\qquad
D_\infty=S_\infty^{-1}B_\infty S_\infty
=\diag(-1/2,0,1/2)
\end{equation}
lautet das System in $t=1/x$:
\begin{equation}
\frac{D_\infty}{t}+\frac{C_\infty}{1-t},
\qquad C_\infty=-S_\infty^{-1}A_0S_\infty,\qquad
(C_\infty)_{31}=0.
\end{equation}
Für eine lokale Lösung $H(t)t^{D_\infty}$, $H_0=I$, ist die Rekursion
\begin{equation}
(n-\operatorname{ad}D_\infty)H_n
=C_\infty\sum_{j=0}^{n-1}H_j.
\end{equation}
Nur $n=1$, Eintrag $(3,1)$, ist resonant. Genau dort verschwindet die rechte Seite; bei $n\ge2$ ist der Rekursionsoperator invertierbar. Es entsteht kein logarithmischer Jordananteil. Die drei lokalen Monodromien haben deshalb tatsächlich die Ordnungen $4,3,2$.

Ihr Bild ist ein Quotient der endlichen sphärischen Dreiecksgruppe $\Delta(4,3,2)\cong S_4$. Ein echter Quotient von $S_4$ besitzt kein Element der Ordnung vier. Das Bild ist daher treu; seine dreidimensionale Darstellung ist irreduzibel. Die Determinante eins wählt die Rotationsdarstellung, also die Standarddarstellung mal Vorzeichencharakter.

Die Markierung stimmt ebenfalls: Aus der Quadratäquivarianz folgt bei $\Psi(0)=I$ die Identität $\Psi(iz)=U\Psi(z)U^{-1}$. Die positive $x=0$-Schleife wirkt deshalb mit $U^{-1}$. Die exakten Zielmatrizen erfüllen $(MU^{-1})^2=I$. Das markierte Paar ist somit analytisch bis auf gemeinsamen Basiswechsel bestimmt. Weil $U$ drei verschiedene Eigenwerte hat, ist ein noch verbleibender Basiswechsel diagonal; insbesondere folgen
\begin{equation}
\boxed{\diag M=(0,-i/2,+i/2),\qquad |M_{22}|=\tfrac12}
\end{equation}
ohne ODE-Numerik. Die genaue gedruckte Offdiagonal-Vorzeichenfassung legt eine Basiszeichenkonvention fest. Der Betrag $1/2$ ist ein Monodromiedatum und nicht der kontinuierliche Transportpol $\delta_{\rm ph}$.

\textbf{Auf der ursprünglichen $z$-Ebene erzeugen die vier Punkturmonodromien $A_4$ mit zwölf Elementen.} Erst die zusätzliche Quadratrotation erweitert das Paket zu $S_4$ mit 24 Elementen. Diese Unterscheidung ist für die Deckfläche wichtig. Die Reduktion, der Resonanzkoeffizient und die Gruppenrelation wurden unabhängig exakt gegengeprüft; der Beweis ist nicht als neuer Lean-Beweis ausgegeben.

Für die exakten Originalmatrizen aus v117 gibt es den unitären Basiswechsel
\begin{equation}
\begin{aligned}
F&=\begin{pmatrix}
1&0&0\\
0&-(1+i)/2&(-1+i)/2\\
0&(1-i)/2&(1+i)/2
\end{pmatrix},\\
F^\dagger MF&=P_c=\begin{pmatrix}0&0&1\\1&0&0\\0&1&0\end{pmatrix},
\qquad
F^\dagger UF=R_c=\begin{pmatrix}1&0&0\\0&0&1\\0&-1&0\end{pmatrix}.
\end{aligned}
\end{equation}
Die vier lokalen Monodromien $U^kMU^{-k}$ sind damit signierte Permutationsmatrizen. Ihre Permutationsanteile sind abwechselnd $P_c,P_c^{-1},P_c,P_c^{-1}$. Genau diese Verzweigung an $z=1,i,-1,-i$ realisiert die dreiblättrige Kurve
\begin{equation}\label{eq:flavorg2}
w^3=\frac{z^2-1}{z^2+1},
\qquad
\eta=z(1-w^3)
\quad\Longrightarrow\quad
\boxed{\eta^2=1-w^6.}
\end{equation}
Die Riemann--Hurwitz-Rechnung ist $2g-2=3(-2)+4(3-1)=2$, also $g=2$. Die sechs hyperelliptischen Verzweigungspunkte sind genau $\mu_6$. Die verbleibenden Vorzeichen der Monodromiematrizen sind Daten eines Linienlokalsystems auf dieser Kurve. Die Kurve allein enthält sie noch nicht.

Die ursprüngliche Quadratrotation hebt ausdrücklich zu
\begin{equation}
r:(w,\eta)\longmapsto
\left(\frac1w,-\frac{i\eta}{w^3}\right),\qquad
r^2(w,\eta)=(w,-\eta),\quad r^4=I.
\end{equation}
Auch der Sechserzyklus ist auf derselben Kurve vorhanden. Für $\omega=e^{2\pi i/3}$ setze
\begin{equation}
c(w,\eta)=(\omega w,\eta),\qquad
s(w,\eta)=(-w,\eta).
\end{equation}
Die sechs Punkte lassen sich als $p_a=\omega^a$, $q_a=-\omega^a$ ordnen. In diesem Punktlabelraum wirken
\begin{equation}
c=\diag(P_c,P_c),\qquad
s=\begin{pmatrix}0&I\\I&0\end{pmatrix},\qquad T=sc.
\end{equation}
Der explizite unitäre Intertwiner
\begin{equation}
J_{\rm pts}=\frac1{\sqrt2}
\begin{pmatrix}F^\dagger&F^\dagger\\F^\dagger&-F^\dagger\end{pmatrix}
\end{equation}
erfüllt
\begin{equation}\label{eq:pointhexagon}
\boxed{J_{\rm pts}^\dagger TJ_{\rm pts}=M\oplus(-M),\qquad
J_{\rm pts}^\dagger sJ_{\rm pts}=\diag(I,-I).}
\end{equation}
Damit ist der endliche Sechseroperator ein tatsächlicher markierter Punkttransport. Er wurde nicht aus der Zahl sechs erraten.

\begin{center}
\begin{tikzpicture}[node distance=7mm]
\node[box,text width=12.8cm] (q) {\textbf{Originales Quadrat}\quad $1,i,-1,-i$\\drei Familien; signierte lokale Monodromie};
\node[box,below=of q,text width=12.8cm] (g) {\textbf{Dieselbe Deckkurve}\quad $w^3=(z^2-1)/(z^2+1)$\\zwei Henkel; Vorzeichen-Linienlokalsystem mitführen};
\node[box,below=of g,text width=12.8cm] (h) {\textbf{Zweite Ansicht}\quad $\eta^2=1-w^6$\\sechs Verzweigungspunkte; $M\oplus(-M)$ als Punkttransport};
\node[box,below=of h,text width=12.8cm] (f) {\textbf{Noch zu identifizierende Feldwirkung}\\derselbe Transport auf geladenen Defekten, ihrem Zustand und den vollständigen Antworten};
\draw[arr] (q)--(g);\draw[arr] (g)--(h);\draw[dasharr] (h)--(f);
\end{tikzpicture}
\end{center}

\picturetext{Vier Marken, drei Blätter und sechs ausgezeichnete Punkte können Ansichten derselben geometrischen Konstruktion sein. Die mitgeführten Vorzeichen sind dabei wie zusätzliche Anweisungen auf den Blättern: Das Drehen der Zeichnung genügt noch nicht, um die Feldoperation richtig zu drehen.}

Insbesondere ist $s$ nicht die hyperelliptische Involution. Letztere fixiert alle sechs Verzweigungspunkte; $s$ vertauscht die beiden Dreierbahnen und induziert $z\mapsto1/z$. Der hier bewiesene Punktintertwiner darf deshalb weder als Feldintertwiner noch als Ableitung des ursprünglichen Nahtprojektors ausgegeben werden.

\subsection{Der Sechsertransport entsteht aus zwei ursprünglichen Halbwegen}\label{sec:pathclock}
\textbf{Ein globaler Inversionsautomorphismus des einzelnen Flachbündels ist hierfür nicht nötig.}
Wähle einen regulären Ausgangspunkt $x\in(0,1)$ und setze $x'=1/x$.
Der Kreis mit Mittelpunkt $a=(x+x')/2$ und Radius $r=(x'-x)/2$ erfüllt
$a^2-r^2=1$ und ist unter $\iota(z)=1/z$ invariant. Sein unterer Halbkreis $p$ führt von $x$ nach $x'$. Der Bildweg $\iota p$ führt auf dem oberen Halbkreis zurück. Zusammen umlaufen sie $+1$ genau einmal positiv. Der ganze Kreis liegt im rechten Halbplan und umschließt keine weitere Fuchs-Punktur.

Seien $P=\operatorname{PT}(p):E_x\to E_{x'}$ und
$Q=\operatorname{PT}(\iota p):E_{x'}\to E_x$ die tatsächlichen Paralleltransporte der ursprünglichen Verbindung. Dann gilt
\begin{equation}
QP=M_x,\qquad PQ=PM_xP^{-1}.
\end{equation}
Die endliche Monodromie besitzt eine positive invariante hermitesche Form, etwa durch Mittelung einer positiven Form über ihre endliche Gruppe. Paralleltransport setzt sie zur flachen Fasermetrik fort. Wählt man den Rahmen bei $x'$ durch Transport des orthonormalen Rahmens bei $x$, so erhält man
\begin{equation}\label{eq:pathT}
P=I,\quad Q=M,\qquad
T=\begin{pmatrix}0&M\\I&0\end{pmatrix},\qquad
T^2=\diag(M,M),\qquad T^6=I.
\end{equation}
Die rohen ODE-Matrizen in einem beliebigen holomorphen Rahmen sind deshalb nicht automatisch bezüglich der gewöhnlichen euklidischen Matrixmetrik unitär. Der Satz verwendet die flache hermitesche Metrik.

Mit dem expliziten unitären Wörterbuch
\begin{equation}
J=\frac1{\sqrt2}\begin{pmatrix}I&I\\M&-M\end{pmatrix}
\end{equation}
folgt exakt
\begin{equation}\label{eq:pathclockdictionary}
\boxed{J^\dagger T^{-1}J=M\oplus(-M)=U_6,\qquad
J^\dagger T^3J=\diag(I,-I).}
\end{equation}
Die ursprüngliche Clock entspricht bei dieser positiven Umlaufkonvention dem inversen Halbschritt. Diese Orientierung wird ausdrücklich mitgeführt. $T^3$ ist eine komplexlineare Involution mit zwei komplex dreidimensionalen Eigenräumen; sie ist kein antiunitärer Reflexionsswap.

Für den bereits gegebenen Flavoroperator folgt damit
\begin{equation}
D^{\rm path}_{y,\delta}=yI-\delta T^{-1},\qquad
\boxed{J^\dagger D^{\rm path}_{y,\delta}J=yI-\delta U_6.}
\end{equation}
Dasselbe $J$ transportiert die positive Gramform, ihre regulierte Resolvente, den inversen Quadratwurzelkernel und die Determinante
$\det D=y^6-\delta^6$. Es werden weder neue Flavorparameter eingesetzt noch bestehende Massen- und Yukawa-Herleitungen ersetzt. Veränderungen von $x$ innerhalb $(0,1)$ ändern nur den parallel transportierten Rahmen und damit den Operator durch unitäre Konjugation.

\begin{center}
\begin{tikzpicture}
\node[box,text width=4.1cm] (left) at (0,0) {Zugang $x$\\drei Familienkomponenten};
\node[box,text width=4.1cm] (right) at (8,0) {Zugang $1/x$\\dieselben Feldtypen};
\draw[arr,bend right=28] (left.south) to node[below,font=\small] {unterer Halbweg: $P$} (right.south);
\draw[arr,bend right=28] (right.north) to node[above,font=\small] {oberer Halbweg: $Q$} (left.north);
\node[align=center,font=\small,text=navy] at (4,-2) {zwei Halbschritte: $M$\qquad sechs Halbschritte: $I$\\Original-Flavoroperator: $yI-\delta T^{-1}$};
\end{tikzpicture}
\end{center}

\picturetext{Ein Paket mit drei Familienkomponenten wird über zwei Zugänge weitergereicht. Nach der Hin- und Rückrunde trägt es die ursprüngliche Familienänderung $M$. Drei solcher Runden schließen sich. Die sechs Schritte sind damit eine Operation am vorhandenen Paket, keine neue Zahlenspielerei.}

\textbf{Der genaue Feldumfang bleibt erhalten.}
$E_x\oplus E_{x'}$ ist ein Raum von Randdaten an zwei Zugängen. Er ist nicht der sechsdimensionale Raum horizontaler Lösungen: Auf einer eingeschnittenen verbundenen Domäne erfüllen deren Werte bereits $s(x')=Ps(x)$ und bilden nur einen dreidimensionalen Graphen. Die Quantisierung der gesamten Feldquelle muss daher den ursprünglichen Testfunktionenraum und seinen Riemann--Hilbert-Kernel mitführen. Die endliche Gleichung~\eqref{eq:pathclockdictionary} allein behauptet weder eine globale Zustandsinvarianz noch die Gleichheit $T^3=\mathsf iC$ mit dem ursprünglichen Calder\'onoperator.

